\par\Needspace{4\baselineskip}
\begin{center}
\hypertarget{native-043-article027}{}\pdfbookmark[2]{27.}{article-027}27.
\end{center}
Differenziert man Gleichung IV und behandelt dabei $u$, $v$ und $\psi$ zugleich als veränderlich, so erhält man
\[
\begin{aligned}
\frac{\mathrm{d}u}{u}&=\frac{\sin\psi\,\mathrm{d}v+\sin v\,\mathrm{d}\psi}{2\cos\frac{1}{2}(v-\psi)\cos\frac{1}{2}(v+\psi)}\\
&=\frac{r\operatorname{tang}\psi}{p}\,\mathrm{d}v+\frac{r\sin v}{p\cos\psi}\,\mathrm{d}\psi.
\end{aligned}
\]

Differenziert man ebenso Gleichung XI, so ergibt sich die folgende Beziehung zwischen den differenziellen Änderungen der Größen $u$, $\psi$ und $N$
\SourceLeaf{044}
\[
\frac{\mathrm{d}N}{\lambda}=\left\{\frac{1}{2}e\left(1+\frac{1}{uu}\right)-\frac{1}{u}\right\}\mathrm{d}u+\frac{(uu-1)\sin\psi}{2u\cos\psi^2}\,\mathrm{d}\psi
\]
oder
\[
\frac{\mathrm{d}N}{\lambda}=\frac{r}{bu}\,\mathrm{d}u+\frac{r\sin v}{b\cos\psi}\,\mathrm{d}\psi.
\]
Eliminiert man daraus $\mathrm{d}u$ mithilfe der vorangehenden Gleichung, so erhält man
\[
\frac{\mathrm{d}N}{\lambda}=\frac{rr}{bb\operatorname{tang}\psi}\,\mathrm{d}v+\left(1+\frac{r}{p}\right)\frac{r\sin v}{b\cos\psi}\,\mathrm{d}\psi
\]
oder
\[
\begin{aligned}
\mathrm{d}v&=\frac{bb\operatorname{tang}\psi}{\lambda rr}\,\mathrm{d}N-\left(\frac{b}{r}+\frac{b}{p}\right)\frac{\sin v\operatorname{tang}\psi}{\cos\psi}\raisebox{-.6ex}{.}\,\mathrm{d}\psi\\
&=\frac{bb\operatorname{tang}\psi}{\lambda rr}\,\mathrm{d}N-\left(1+\frac{p}{r}\right)\frac{\sin v}{\sin\psi}\,\mathrm{d}\psi.
\end{aligned}
\]

\par\Needspace{4\baselineskip}
\begin{center}
\hypertarget{native-044-article028}{}\pdfbookmark[2]{28.}{article-028}28.
\end{center}
Wir differenzieren Gleichung X, wobei wir $r$, $b$, $e$ und $u$ sämtlich als veränderlich behandeln. Setzen wir $\mathrm{d}e=\frac{\sin\psi}{\cos\psi^2}\,\mathrm{d}\psi$ ein und eliminieren wir $\mathrm{d}u$ mithilfe der im vorigen Artikel angegebenen Gleichung zwischen $\mathrm{d}N$, $\mathrm{d}u$ und $\mathrm{d}\psi$, so erhalten wir
\[
\begin{aligned}
\mathrm{d}r&=\frac{r}{b}\,\mathrm{d}b+\frac{bbe(uu-1)}{2\lambda ur}\,\mathrm{d}N\\
&\quad+\frac{b}{2\cos\psi^2}\left\{\left(u+\frac{1}{u}\right)\sin\psi-\left(u-\frac{1}{u}\right)\sin v\right\}\mathrm{d}\psi.
\end{aligned}
\]
Nach Gleichung VIII wird der Koeffizient von $\mathrm{d}N$ zu $\frac{b\sin v}{\lambda\sin\psi}$. Setzt man im Koeffizienten von $\mathrm{d}\psi$ nach Gleichung IV $u(\sin\psi-\sin v)=\sin(\psi-v)$ und $\frac{1}{u}(\sin\psi+\sin v)=\sin(\psi+v)$ ein, so wird er zu $\frac{b\sin\psi\cos v}{\cos\psi^2}=\frac{p\cos v}{\sin\psi}$. Damit erhält man
\[
\mathrm{d}r=\frac{r}{b}\,\mathrm{d}b+\frac{b\sin v}{\lambda\sin\psi}\,\mathrm{d}N+\frac{p\cos v}{\sin\psi}\,\mathrm{d}\psi.
\]
Betrachtet man ferner $N$ als Funktion von $b$ und $t$, so erhält man $\mathrm{d}N=\frac{N}{t}\,\mathrm{d}t-\frac{3}{2}\cdot\frac{N}{b}\,\mathrm{d}b$. Setzt man diesen Wert ein, so erhält man $\mathrm{d}r$ und ebenso im vorigen Artikel $\mathrm{d}v$ als Ausdrücke in $\mathrm{d}t$, $\mathrm{d}b$ und $\mathrm{d}\psi$. Auch hier ist die obige Bemerkung zu wiederholen. Will man die Änderungen der Winkel $v$ und $\psi$ von Teilen des Radius in Bogensekunden umrechnen, so sind entweder alle Terme, die $\mathrm{d}v$ oder $\mathrm{d}\psi$ enthalten, durch $206264{,}8$ zu teilen oder alle übrigen Terme mit dieser Zahl zu multiplizieren.
